\section{Анализ результатов}

\subsection{Ключевые числа}
\begin{table}[H]
\centering
\begin{tabular}{lccc}
\toprule
Аппроксимация & RMSE & $L_\infty$ & RelRMSE \\
\midrule
full\_quadratic & 1.373e-05 & 2.814e-05 & 1.372e-02 \\
diagonal\_quadratic & 1.539e-05 & 3.401e-05 & 1.538e-02 \\
hessian\_quadratic & 1.194e-04 & 3.355e-04 & 1.194e-01 \\
\bottomrule
\end{tabular}
\caption{Метрики для model1 при $r=0.1$.}
\label{tab:metrics_r01}
\end{table}

\subsection{Интерпретация}
\begin{itemize}
    \item $q_1$ показывает наилучшее качество и подтверждает гипотезу о локальной квадратичности.
    \item $q_2$ близка к $q_1$, что говорит о малом вкладе смешанного члена в центре.
    \item $q_3$ заметно хуже: FD-оценка чувствительна к численному шуму и выбору шага $\epsilon$.
\end{itemize}

\subsection{Почему это научно важно}
Результат демонстрирует не просто «построение красивой картинки», а проверку гипотезы о структуре ландшафта с количественной валидацией.
